% answers_short.tex —— 2017 年·天津卷·理科 · 简短答案速查
% 由 tools/make_short_answers.py 自 answers.tex 的 \gksolhead/\gkanswer 提取，
% 并与 2017/2017_数学真题_天津理_答案要点.md 逐题比对后生成，勿手改（改源 MD 后重跑）。

\begin{gknotice}
\textbf{使用说明}：本节只给最终答案便于快速核对，完整推导见解析版。答案为编写者独立推导，\textbf{非官方标准答案}。
\end{gknotice}

\gkpart{选择题}
\begin{center}\begin{tabular}{ccccccccc}
\toprule
题号 & 1 & 2 & 3 & 4 & 5 & 6 & 7 & 8 \\
\midrule
答案 & B & D & C & A & B & C & A & A \\
\bottomrule
\end{tabular}\end{center}

\gkpart{填空题}
\begin{itemize}[leftmargin=2.6em]
  \item[\textbf{9.}] $-2$
  \item[\textbf{10.}] $9\pi/2$
  \item[\textbf{11.}] $2$
  \item[\textbf{12.}] 最小值 $4$
  \item[\textbf{13.}] $3/11$
  \item[\textbf{14.}] $1080$
\end{itemize}

\gkpart{解答题}
\begin{itemize}[leftmargin=2.6em]
  \item[\textbf{15.}] (1) $b=\sqrt{13}$，$\sin A=3/\sqrt{13}$；(2) $\sin(2A+\pi/4)=7\sqrt2/26$
  \item[\textbf{16.}] (1) $P(X=0,1,2,3)=(1/4,11/24,1/4,1/24)$，$E(X)=13/12$；(2) $11/48$
  \item[\textbf{17.}] (1) $MN\parallel$ 平面 $BDE$；(2) 两平面夹角的正弦为 $\sqrt{5/21}$；(3) $AH=1/2$ 或 $8/5$
  \item[\textbf{18.}] (1) $a_n=3n-2,b_n=2^n$；(2) $\sum_{k=1}^n a_{2k}b_{2k-1}=[8+(3n-2)4^{n+1}]/3$
  \item[\textbf{19.}] (1) 椭圆 $x^2+y^2/(3/4)=1$，抛物线 $y^2=4x$；(2) 直线 $y=\pm(\sqrt6/2)(x-1)$
  \item[\textbf{20.}] (1) $g$ 在 $(-\infty,-1)$、$(1/4,+\infty)$ 递增，在 $(-1,1/4)$ 递减；(2) 见解析版；(3) 见解析版
\end{itemize}
